%==============================================================================
% SAMPLE — Chapter 2, Sections 2.3 and 2.4, for the rebuilt report.
% Drop-in: uses the existing preamble macros (\Kst, \Zst, \U, \rset, \confs)
% and the existing theorem environments.
%==============================================================================

\subsection{Grouping the jobs}\label{ssec:jgp}

Section~\ref{ssec:ktns} settled the loading problem. What remains is the order.
Solving for the order directly is hard, so practice does something indirect. It
first asks a simpler question: which jobs could share a magazine?

\begin{definition}[Feasible group]\label{def:group}
A set of jobs $G \subseteq J$ is a \emph{feasible group} if the tools its jobs
need, taken together, fit in the magazine:
\[
  \Bigl\lvert \bigcup_{j \in G} T_j \Bigr\rvert \;\le\; b .
\]
\end{definition}

A feasible group can be run start to finish without a single tool change. Load
the tools its jobs need, run them in any order, and only then change anything.
This makes the following question natural.

\begin{definition}[Job Grouping Problem]\label{def:jgp}
The \emph{Job Grouping Problem} (JGP) is the problem of partitioning $J$ into as
few feasible groups as possible. The smallest possible number of groups is
written $\Kst$.
\end{definition}

The JGP ignores order completely. It asks only how the jobs may be divided. It is
itself NP-hard, but it is much easier in practice than the SSP, and it can be
solved exactly on the instance sizes considered here.

Two problems now sit next to each other. The JGP says which jobs may share a
magazine. The SSP asks for an order. The standard heuristic joins them, in three
steps.

\begin{enumerate}[topsep=4pt,itemsep=2pt]
  \item \textbf{Group.} Solve the JGP. This returns $\Kst$ feasible groups
  $G_1,\dots,G_{\Kst}$, and for each group a magazine \emph{configuration}
  $C_l$ of size $b$ containing everything the group needs.
  \item \textbf{Sequence.} Decide the order of the groups. The cost of moving
  from configuration $C$ to configuration $C'$ is the number of tools that must
  be inserted, $d(C,C') = \lvert C' \setminus C \rvert$. Finding the cheapest
  order of the $\Kst$ configurations is a travelling-salesman problem on $\Kst$
  points, which is small enough to solve exactly.
  \item \textbf{Evaluate.} Write out the jobs group by group, in the order chosen
  in step~2, to obtain a single job order. Apply KTNS to that order and report
  its cost.
\end{enumerate}

This is the method studied throughout this report. It is fast, it is what
practitioners use, and its output is always a valid schedule, so its cost is
never below the optimum $\Zst$.

\begin{example}[The heuristic on the $6$-ring]\label{ex:jgp-ring}
Take the $6$-ring of Example~\ref{ex:ring}: six jobs, six tools, $b=3$, and
$T_j = \rset{j, j+1}$ with indices modulo~$6$.

Step~1 returns $\Kst = 3$. One optimal grouping is
\[
  G_1 = \rset{1,2}, \qquad G_2 = \rset{3,4}, \qquad G_3 = \rset{5,6},
\]
whose tool requirements are $\rset{1,2,3}$, $\rset{3,4,5}$ and $\rset{5,6,1}$.
Each has exactly three tools, so each group fixes its configuration with no
freedom left: $C_1 = \rset{1,2,3}$, $C_2 = \rset{3,4,5}$, $C_3 = \rset{5,6,1}$.

Step~2 compares the three configurations. Any two of them share exactly one
tool, so moving between any two costs two insertions. Every order of the three
groups therefore costs the same, and step~2 may return any of them. Take
$G_1, G_2, G_3$.

Step~3 writes out the jobs group by group, giving the order $1,2,3,4,5,6$, and
applies KTNS to it. That is the order already traced in Example~\ref{ex:ring},
and its cost is $6$ insertions, or $3$ once the first fill is not charged.
\end{example}

Two things in this example are worth pausing on, because both matter later.

First, the method does not fully determine its own answer. Step~1 may return
several optimal groupings; step~2 may face ties, as it does here; and the order
of the jobs \emph{inside} a group is left open by all three steps. Different
choices can give different costs. Throughout this report the heuristic is read
in its most favourable form: its cost is the best value reachable over all these
choices. This is the reading that makes the comparison with $\Zst$ meaningful,
because any weaker reading would measure the tie-breaking rule rather than the
method.

Second, and less obviously, step~3 does something that step~2 did not
anticipate. This is the subject of the next section.

\subsection{Two ways to measure the heuristic}\label{ssec:frames}

Steps~1 and~2 reason entirely in terms of configurations. A group is assigned one
magazine, and that magazine is imagined to stay loaded while the whole group
runs. Under that picture the schedule is a walk through $\Kst$ configurations,
and its cost is the sum of the transition costs $d(C,C')$ along the walk.

Step~3 does not do this. It discards the configurations, keeps only the job
order they produced, and hands that order to KTNS. KTNS is free to change the
magazine between any two consecutive jobs, including two jobs of the same group.
The group structure survives only as an ordering of the jobs. It does not
constrain the magazine.

The two readings give different numbers, and both appear in the literature
without being distinguished. This report separates them.

\begin{definition}[The walk value and the heuristic value]\label{def:frames}
Let $\mathcal{P}$ range over the partitions of $J$ into $\Kst$ feasible groups,
and for each such partition let the configurations be chosen and ordered as
cheaply as possible. Write
\[
  H_{\mathrm{walk}}
  \;=\;\min_{\mathcal{P}}\;\min_{\sigma}\;
        \sum_{l=1}^{\Kst - 1} d\bigl(C_{\sigma(l)}, C_{\sigma(l+1)}\bigr)
\]
for the cheapest such walk. Write $H$ for the cost of the two-phase heuristic as
defined in Section~\ref{ssec:jgp}: the minimum, over the same partitions, over
all orders of the groups and all orders of the jobs within each group, of the
KTNS cost of the resulting job order.
\end{definition}

\begin{proposition}\label{prop:frames}
For every instance, $H \le H_{\mathrm{walk}}$.
\end{proposition}

\begin{proof}
Fix a partition and an order of its groups attaining $H_{\mathrm{walk}}$. Writing
the jobs out group by group in that order gives a job order. Holding each
group's configuration fixed while its jobs run is one admissible loading for
that order, and its cost is exactly the walk cost. By
Proposition~\ref{prop:ktns}, KTNS returns the cheapest loading for that order,
so its cost is at most the walk cost. Since $H$ minimises over at least this
one job order, $H \le H_{\mathrm{walk}}$.
\end{proof}

The inequality can be strict, and it is strict on the running example.

\begin{example}[The two values differ on the $6$-ring]\label{ex:frames-ring}
Continue Example~\ref{ex:jgp-ring}. The walk pays three insertions to build
$C_1 = \rset{1,2,3}$, then two to reach $C_2 = \rset{3,4,5}$, then two to reach
$C_3 = \rset{5,6,1}$. That is $H_{\mathrm{walk}} = 7$, or $4$ free-initial.

KTNS on the same job order costs $6$, or $3$ free-initial
(Example~\ref{ex:ring}). Since $\Zst = 3$, the heuristic is \emph{optimal} on the
$6$-ring, while the walk value is one above the optimum.

Figure~\ref{fig:frames} shows where the two part company. The walk insists that
the middle group runs under $C_2 = \rset{3,4,5}$ and nothing else. Tool~1 is not
in $C_2$, so it is evicted at position~3 and must be reloaded at position~5 for
job~6. KTNS is not bound to $C_2$. It holds $\rset{1,3,4}$ at position~3 and
$\rset{1,4,5}$ at position~4 --- two different magazines inside one group, each
serving its own job --- and so keeps tool~1 resident throughout. The one
re-insertion is avoided, and the schedule is optimal.
\end{example}

\begin{figure}[t]\centering
\begin{tikzpicture}[x=0.8cm,y=0.8cm,font=\footnotesize]
 % ---------------- (a) the configuration walk ----------------
 \begin{scope}
  \node[font=\footnotesize\bfseries,anchor=west] at (0.15,7.6)
     {(a) the walk: one magazine per group};
  \foreach \p/\req in {1/{1{,}2},2/{2{,}3},3/{3{,}4},4/{4{,}5},5/{5{,}6},6/{6{,}1}}{
    \node[font=\footnotesize] at (\p,6.05) {\p};
    \node[font=\scriptsize,text=black!55] at (\p,6.7) {$\{\req\}$};}
  \foreach \t/\y in {1/5,2/4,3/3,4/2,5/1,6/0}{
    \node[anchor=east,font=\footnotesize] at (0.4,\y) {$t_{\t}$};}
  \foreach \p/\y in {1/5,2/5,5/5,6/5, 1/4,2/4, 1/3,2/3,3/3,4/3,
                     3/2,4/2, 3/1,4/1,5/1,6/1, 5/0,6/0}{
     \fill[cBlue!20] (\p-0.5,\y-0.5) rectangle (\p+0.5,\y+0.5);}
  \foreach \p/\y in {1/5,1/4,1/3, 3/2,3/1, 5/5,5/0}{
     \fill[cOrange!90] (\p-0.5,\y-0.5) rectangle (\p+0.5,\y+0.5);
     \draw[white,line width=0.9pt] (\p-0.17,\y)--(\p+0.17,\y);
     \draw[white,line width=0.9pt] (\p,\y-0.17)--(\p,\y+0.17);}
  \draw[black!35,step=1] (0.5,-0.5) grid (6.5,5.5);
  \draw[black!55,semithick,dashed] (2.5,-0.5) -- (2.5,5.5);
  \draw[black!55,semithick,dashed] (4.5,-0.5) -- (4.5,5.5);
  \foreach \x/\g in {1.5/1,3.5/2,5.5/3}{
    \node[font=\scriptsize,text=black!55] at (\x,-1.05) {$G_{\g}$};}
  \node[font=\scriptsize,text=black!70,anchor=west] at (0.5,-1.8)
     {seven insertions};
 \end{scope}
 % ---------------- (b) KTNS on the same job order ----------------
 \begin{scope}[xshift=7.7cm]
  \node[font=\footnotesize\bfseries,anchor=west] at (0.15,7.6)
     {(b) KTNS on the same job order};
  \foreach \p/\req in {1/{1{,}2},2/{2{,}3},3/{3{,}4},4/{4{,}5},5/{5{,}6},6/{6{,}1}}{
    \node[font=\footnotesize] at (\p,6.05) {\p};
    \node[font=\scriptsize,text=black!55] at (\p,6.7) {$\{\req\}$};}
  \foreach \t/\y in {1/5,2/4,3/3,4/2,5/1,6/0}{
    \node[anchor=east,font=\footnotesize] at (0.4,\y) {$t_{\t}$};}
  \foreach \p/\y in {1/5,2/5,3/5,4/5,5/5,6/5, 1/4,2/4, 1/3,2/3,3/3,
                     3/2,4/2, 4/1,5/1,6/1, 5/0,6/0}{
     \fill[cBlue!20] (\p-0.5,\y-0.5) rectangle (\p+0.5,\y+0.5);}
  \foreach \p/\y in {1/5,1/4,1/3, 3/2, 4/1, 5/0}{
     \fill[cOrange!90] (\p-0.5,\y-0.5) rectangle (\p+0.5,\y+0.5);
     \draw[white,line width=0.9pt] (\p-0.17,\y)--(\p+0.17,\y);
     \draw[white,line width=0.9pt] (\p,\y-0.17)--(\p,\y+0.17);}
  \draw[black!35,step=1] (0.5,-0.5) grid (6.5,5.5);
  \draw[black!55,semithick,dashed] (2.5,-0.5) -- (2.5,5.5);
  \draw[black!55,semithick,dashed] (4.5,-0.5) -- (4.5,5.5);
  \foreach \x/\g in {1.5/1,3.5/2,5.5/3}{
    \node[font=\scriptsize,text=black!55] at (\x,-1.05) {$G_{\g}$};}
  \node[font=\scriptsize,text=black!70,anchor=west] at (0.5,-1.8)
     {six insertions};
 \end{scope}
\end{tikzpicture}
\caption{The same grouping of the $6$-ring, measured two ways. Columns are the
six processing positions, with the tools each job needs above them; rows are the
six tools. A shaded cell means the tool sits in the magazine at that position;
a ``$+$'' marks an \emph{insertion}. The dashed lines separate the three groups.
In (a) each group runs under one fixed magazine, so tool~$1$ (top row) is evicted
at position~$3$ and reloaded at position~$5$ for job~$6$: seven insertions. In
(b) KTNS may change the magazine inside a group, so tool~$1$ is loaded once and
kept across all six positions, and the count is six. That single re-insertion is
the whole difference between $H_{\mathrm{walk}}$ and $H$.}
\label{fig:frames}
\end{figure}

\begin{remark}[Why the distinction is kept throughout]\label{rem:frames}
It would be simpler to work with $H_{\mathrm{walk}}$ alone. It has a clean
description --- a shortest walk through $\Kst$ configurations --- and the
structural results of Chapter~\ref{sec:gap} are naturally proved about it.

That simplification is not available, for two reasons. The first is that
$H_{\mathrm{walk}}$ is not what the method returns. A practitioner running the
three steps of Section~\ref{ssec:jgp} obtains $H$, and any statement about how
far the method can be from optimal must be a statement about $H$. The second is
that the two can differ on exactly the instances that matter: on the $6$-ring,
the family this report returns to most often, $H_{\mathrm{walk}}$ exceeds the
optimum while $H$ attains it.

Every result in Chapter~\ref{sec:gap} is therefore stated with its value named.
Proposition~\ref{prop:frames} carries upper bounds from one to the other:
anything proved to bound $H_{\mathrm{walk}}$ from above also bounds $H$. Nothing
transfers in the other direction, and in particular an instance on which
$H_{\mathrm{walk}}$ exceeds $\Zst$ is not, by itself, an instance on which the
heuristic is suboptimal.
\end{remark}
